\section{Experiments and Results}
\label{sec:results}

\paragraph{Evaluation protocol.}
One match yields many team-phase rows, so a row-level split would leak the
same match context into train and test. We therefore split by match, using
17,460 train, 2,182 validation, and 2,182 test matches with seed 42. The main
comparison uses 134,245 aligned risk-set cells from the test matches. Each
cell is one \((match, team, phase)\) observation for a team that is alive at
\(b_k\) and has an observed next phase. Paired significance tests aggregate
performance at the match level, so that a match with many teams or phases
does not dominate the test.
We define the reported metrics as follows. AUC is the area under the ROC curve and summarizes how
well predicted hazards rank eliminated cells above non-eliminated cells; higher
is better. AP is average precision, the area under the precision-recall curve,
and complements AUC when positive hazard cells are less frequent; higher is
better. BCE is binary cross-entropy between the hazard label and predicted
probability. The Brier score is the mean squared error of predicted
probabilities. ECE is expected calibration error: we split predicted
probabilities into 10 equal-width bins and average, weighted by bin size, the
gap between mean predicted hazard and observed elimination rate. BCE, Brier,
and ECE are lower-is-better.

\paragraph{Models.}
We compare a histogram gradient boosting (HGB) feature-only baseline, a
no-edge DynamicSTGNN that keeps the temporal architecture but zeros all
adjacency channels, and three graph variants: proximity only, proximity plus
combat, and the full graph with proximity, combat, and co-pressure. HGB is a
strong tabular baseline; no-edge DynamicSTGNN controls for neural
architecture and temporal modeling.

\begin{table}[htbp]
\centering
\caption{Aligned test-set hazard prediction results on 134,245 cells.}
\label{tab:main}
\small
\begin{tabular}{lrrrr}
\toprule
Model & AUC & AP & BCE & ECE \\
\midrule
HGB features & 0.7651 & 0.6626 & 0.5140 & 0.0075 \\
No graph & 0.7681 & 0.6695 & 0.5090 & 0.0039 \\
Proximity & 0.7741 & 0.6762 & 0.5046 & 0.0023 \\
Prox.+combat & 0.7755 & 0.6779 & 0.5036 & 0.0028 \\
Full graph & \textbf{0.7782} & \textbf{0.6808} & \textbf{0.5016} & \textbf{0.0015} \\
\bottomrule
\end{tabular}
\end{table}

Table~\ref{tab:main} shows two patterns. First, team-phase features already
make the hazard learnable: HGB reaches AUC 0.765 and the no-edge neural
control reaches 0.768. Second, graph relations add a smaller but consistent
gain. The full graph improves over no graph by +0.010 AUC and +0.011 AP, and
over HGB by +0.013 AUC and +0.018 AP. Match-paired Wilcoxon tests are
significant for full versus no graph (AUC \(p=1.17\times10^{-30}\), AP
\(p=3.12\times10^{-23}\)) and full versus HGB (AUC
\(p=2.60\times10^{-37}\), AP \(p=9.28\times10^{-36}\)). The full model also
has the best Brier score (0.1662) and ECE (0.0015), with mean predicted
hazard close to the empirical positive rate.

\begin{table}[htbp]
\centering
\caption{Auxiliary squad map-domain experiments. Each map is trained and split
separately; full graph is compared with no-graph DynamicSTGNN.}
\label{tab:map-domain}
\small
\begin{tabular}{lrrrr}
\toprule
Map & Test & No graph & Full & \(\Delta\)AUC \\
\midrule
Baltic & 606 & 0.7653 & 0.7759 & +0.0106 \\
Desert & 149 & 0.7395 & 0.7458 & +0.0063 \\
Neon & 102 & 0.7165 & 0.7244 & +0.0079 \\
Savage & 90 & 0.7338 & 0.7457 & +0.0119 \\
Tiger & 405 & 0.7546 & 0.7605 & +0.0059 \\
\bottomrule
\end{tabular}
\end{table}

Table~\ref{tab:map-domain} shows that the relational gain has the same
direction in separate squad map-domain runs. Full graph improves AUC on all
five major squad maps and AP also increases in every map (+0.003--+0.020).
This auxiliary check is restricted to squad maps with enough data.

\begin{table}[htbp]
\centering
\caption{Representation ablation with the same HGB pipeline. STDB
spatio-temporal features include trajectory and zone/enemy exposure features.}
\label{tab:repr}
\small
\begin{tabular}{lrrr}
\toprule
Feature set & \(d\) & AUC & AP \\
\midrule
Prior-style score summary & 5 & 0.6490 & 0.4547 \\
Score raw inputs & 12 & 0.7184 & 0.5703 \\
STDB spatio-temporal only & 16 & 0.7124 & 0.5895 \\
Conventional aggregates & 57 & 0.7522 & 0.6476 \\
Conventional + STDB spatio-temporal & 72 & \textbf{0.7652} & \textbf{0.6626} \\
\bottomrule
\end{tabular}
\vspace{2pt}
\begin{minipage}{0.96\linewidth}
\scriptsize
\textit{Note.} Score raw inputs are 12 variables covering phase index, combat
counts, damage, knockdowns, distinct opponents damaged, path length, average
speed, 100m enemy density, item pickups, unique items, and care-package
pickups. The score summary uses phase index plus fight score, active score,
strategy type, and type survival rate derived from those raw inputs.
\end{minipage}
\end{table}

Table~\ref{tab:repr} separates representation from model class. Restricting
the learner to the five-variable score summary reduces AUC to 0.649. Adding
the score raw inputs, conventional telemetry aggregates, and STDB trajectory/exposure
features raises AUC to 0.765. Thus the added spatio-temporal features provide
signal beyond conventional aggregates, and the dominant gain is the
telemetry-based phase representation itself, which raises pooled AUC from
0.649 to 0.765 over the score-summary baseline. The dynamic graph supplies an additional relational signal. A graph
shuffle control supports this interpretation: shuffling team identities within
observed graphs at inference drops the full model to AUC 0.701 over five
repeats, below the no-graph control.
